\chapter{The relational program logic}

This chapter develops the coupling-based probabilistic relational Hoare logic
(pRHL) that underlies CatCrypt's game-hopping proofs. The semantics rests on
\emph{probabilistic couplings} of sub-distributions; the judgment
$\vDash \couple{\Phi} c_1 \sim c_2 \couple{\Psi}$ asserts the existence of a
coupling of the two output distributions whose support respects the
postcondition. The structural rules are proved sound directly against this
semantics, and the advantage-level corollaries connect the logic to the
distinguishing advantage used in security reductions.

\section{Couplings and lifting}

\begin{lemma}[Monotonicity of lifting]
  \label{lem:liftr-mono}
  \uses{def:liftr}
  \lean{CatCrypt.Prob.liftR_mono}
  \leanok
  If $R \subseteq S$ pointwise (i.e.\ $R\,a\,b \to S\,a\,b$ for all $a,b$) and
  $d_1 \liftrel{R} d_2$, then $d_1 \liftrel{S} d_2$.
\end{lemma}

\begin{lemma}[Symmetry of lifting]
  \label{lem:liftr-symm}
  \uses{def:liftr}
  \lean{CatCrypt.Prob.liftR_symm}
  \leanok
  If $d_1 \liftrel{R} d_2$, then $d_2 \liftrel{R^{-1}} d_1$, where
  $R^{-1}\,b\,a := R\,a\,b$; the witnessing coupling is obtained by swapping the
  two coordinates of the joint distribution.
\end{lemma}

\begin{lemma}[Equality coupling forces equality]
  \label{lem:liftr-eq-implies-eq}
  \uses{def:liftr}
  \lean{CatCrypt.Prob.liftR_eq_implies_eq}
  \leanok
  If $d_1 \liftrel{(=)} d_2$, then $d_1 = d_2$. An equality coupling places all
  mass on the diagonal, so both marginal sums collapse to the same diagonal
  values.
\end{lemma}

\section{Relational assertions and the judgment}

\begin{definition}[Relational precondition]
  \label{def:rpre}
  \lean{CatCrypt.Relational.RPre}
  \leanok
  A \emph{relational precondition} $\RPre$ is a predicate on pairs of heaps,
  $\Heap \to \Heap \to \Prop$.
\end{definition}

\begin{definition}[Relational postcondition]
  \label{def:rpost}
  \lean{CatCrypt.Relational.RPost}
  \leanok
  A \emph{relational postcondition} $\RPost\,\alpha\,\beta$ is a predicate
  relating two result--heap pairs,
  $\alpha \to \Heap \to \beta \to \Heap \to \Prop$.
\end{definition}

\begin{definition}[Equality precondition]
  \label{def:eqpre}
  \uses{def:rpre}
  \lean{CatCrypt.Relational.eqPre}
  \leanok
  The equality precondition $\eqPre$ holds of $(h_1,h_2)$ exactly when
  $h_1 = h_2$.
\end{definition}

\begin{definition}[Equality postcondition]
  \label{def:eqpost}
  \uses{def:rpost}
  \lean{CatCrypt.Relational.eqPost}
  \leanok
  The equality postcondition $\eqPost$ relates $(a_1,h_1)$ and $(a_2,h_2)$ when
  both the results and the heaps agree, $a_1 = a_2 \wedge h_1 = h_2$.
\end{definition}

\begin{definition}[The pRHL judgment]
  \label{def:rhoare}
  \uses{def:spcomp,def:liftr,def:rpre,def:rpost}
  \lean{CatCrypt.Relational.rHoare}
  \leanok
  For programs $c_1 : \SPComp\,\alpha$, $c_2 : \SPComp\,\beta$, a precondition
  $\Phi : \RPre$ and a postcondition $\Psi : \RPost\,\alpha\,\beta$, the
  judgment $\rHoare\,\Phi\,c_1\,c_2\,\Psi$, written
  $\vDash \couple{\Phi} c_1 \sim c_2 \couple{\Psi}$, is defined as: for all heaps
  $h_1,h_2$ with $\Phi\,h_1\,h_2$, the output distributions are coupled along the
  postcondition,
  \[
    (c_1\,h_1) \;\liftrel{\,\lambda p_1\,p_2.\ \Psi\,p_1.1\,p_1.2\,p_2.1\,p_2.2\,}\; (c_2\,h_2).
  \]
\end{definition}

\section{Structural rules}

\begin{theorem}[Consequence]
  \label{thm:rhoare-conseq}
  \uses{def:rhoare,lem:liftr-mono}
  \lean{CatCrypt.Relational.rHoare_conseq}
  \leanok
  If $\Phi' \Rightarrow \Phi$ (pointwise), $\Psi \Rightarrow \Psi'$ (pointwise),
  and $\rHoare\,\Phi\,c_1\,c_2\,\Psi$, then $\rHoare\,\Phi'\,c_1\,c_2\,\Psi'$.
\end{theorem}

\begin{theorem}[Symmetry]
  \label{thm:rhoare-symm}
  \uses{def:rhoare,lem:liftr-symm}
  \lean{CatCrypt.Relational.rHoare_symm}
  \leanok
  The judgment is symmetric: from
  $\rHoare\,(\lambda h_1\,h_2.\ \Phi\,h_2\,h_1)\,c_2\,c_1\,(\lambda a_2\,h_2\,a_1\,h_1.\ \Psi\,a_1\,h_1\,a_2\,h_2)$
  one derives $\rHoare\,\Phi\,c_1\,c_2\,\Psi$.
\end{theorem}

\begin{theorem}[Sequential composition / bind]
  \label{thm:rhoare-bind}
  \uses{def:rhoare,thm:liftr-bind}
  \lean{CatCrypt.Relational.rHoare_bind}
  \leanok
  Suppose $\rHoare\,\Phi\,c_1\,c_2\,\Psi$ and, for all $a,b$,
  $\rHoare\,(\lambda h_1\,h_2.\ \Psi\,a\,h_1\,b\,h_2)\,(f_1\,a)\,(f_2\,b)\,\Theta$.
  Then
  \[
    \rHoare\,\Phi\,(c_1 \bind f_1)\,(c_2 \bind f_2)\,\Theta .
  \]
\end{theorem}

\begin{theorem}[Reflexivity]
  \label{thm:rhoare-refl}
  \uses{def:rhoare,def:eqpre,def:eqpost,lem:liftr-refl}
  \lean{CatCrypt.Relational.rHoare_refl}
  \leanok
  Every program is related to itself under equality: for all $c$,
  $\rHoare\,\eqPre\,c\,c\,\eqPost$.
\end{theorem}

\section{Sampling rules}

\begin{theorem}[Synchronized sampling, diagonal coupling]
  \label{thm:rhoare-sample-same}
  \uses{def:rhoare,thm:liftr-bind,lem:liftr-refl}
  \lean{CatCrypt.Relational.rHoare_sample_same}
  \leanok
  Sampling the same finite nonempty type on both sides with the diagonal
  coupling yields equal samples:
  \[
    \rHoare\,\Phi\,(\mathsf{sample}\,\alpha)\,(\mathsf{sample}\,\alpha)\,
      (\lambda a_1\,h_1\,a_2\,h_2.\ \Phi\,h_1\,h_2 \wedge a_1 = a_2).
  \]
\end{theorem}

\begin{theorem}[Bijection sampling]
  \label{thm:rhoare-sample-bij}
  \uses{def:rhoare,thm:liftr-bind,lem:liftr-uniform-bij}
  \lean{CatCrypt.Relational.rHoare_sample_bij}
  \leanok
  Sampling along a bijection $f : \alpha \simeq \beta$ relates the two samples by
  $f$:
  \[
    \rHoare\,\Phi\,(\mathsf{sample}\,\alpha)\,(\mathsf{sample}\,\beta)\,
      (\lambda a\,h_1\,b\,h_2.\ \Phi\,h_1\,h_2 \wedge f\,a = b).
  \]
\end{theorem}

\section{Frame rules}

\begin{theorem}[Frame]
  \label{thm:r-frame}
  \uses{def:rhoare,lem:liftr-mono}
  \lean{CatCrypt.Relational.r_frame}
  \leanok
  If $\rHoare\,\Phi\,c_1\,c_2\,\Theta$ and the frame $\Psi$ follows from the
  postcondition ($\Theta\,a\,h_1\,b\,h_2 \to \Psi\,h_1\,h_2$), then the frame may
  be carried alongside:
  \[
    \rHoare\,(\Phi \wedge \Psi)\,c_1\,c_2\,(\lambda a\,h_1\,b\,h_2.\ \Theta\,a\,h_1\,b\,h_2 \wedge \Psi\,h_1\,h_2).
  \]
\end{theorem}

\begin{definition}[Assertion locality]
  \label{def:depends-on}
  \uses{def:rpre}
  \lean{CatCrypt.Relational.DependsOn}
  \leanok
  $\DependsOn\,P\,L$ holds when the assertion $P$ only inspects locations in the
  set $L$: whenever two heap pairs agree on all locations of $L$, $P$ cannot
  distinguish them.
\end{definition}

\begin{definition}[Computation footprint]
  \label{def:preserves-outside}
  \uses{def:spcomp}
  \lean{CatCrypt.Relational.PreservesOutside}
  \leanok
  $\PreservesOutside\,c\,L$ holds when $c$ never modifies locations outside $L$:
  for every output $(a,h')$ in the support of $c\,h$ and every $\mathit{id} \notin L$,
  the value of $\mathit{id}$ in $h'$ equals its value in $h$.
\end{definition}

\begin{theorem}[Location-based frame rule]
  \label{thm:r-frame-local}
  \uses{def:rhoare,def:depends-on,def:preserves-outside}
  \lean{CatCrypt.Relational.r_frame_local}
  \leanok
  Suppose both programs preserve everything outside $L_c$, the frame assertion
  depends only on $L_f$, and $L_c$ and $L_f$ are disjoint. Then, from
  $\rHoare\,\Phi\,c_1\,c_2\,\Psi$,
  \[
    \rHoare\,(\Phi \wedge \mathit{Frame})\,c_1\,c_2\,
      (\lambda a\,h_1\,b\,h_2.\ \Psi\,a\,h_1\,b\,h_2 \wedge \mathit{Frame}\,h_1\,h_2).
  \]
  This is the CatCrypt analogue of the separation-logic frame rule.
\end{theorem}

\section{Reordering and swap rules}

\begin{theorem}[Reorder a pure prefix on the left]
  \label{thm:rhoare-reorder-pure-l}
  \uses{def:rhoare}
  \lean{CatCrypt.Relational.rHoare_reorder_pure_l}
  \leanok
  If $c$ is heap-independent ($\mathsf{IsPure}\,c$), the pure operation may be
  commuted past an independent computation $d$ on the left side:
  from a judgment for $d \bind (\lambda y.\ c \bind (\lambda x.\ k\,x\,y))$ one
  obtains the judgment for $c \bind (\lambda x.\ d \bind (\lambda y.\ k\,x\,y))$
  against the same right side.
\end{theorem}

\begin{theorem}[Swap two samples in opposite order]
  \label{thm:rhoare-sample-comm}
  \uses{def:rhoare,thm:liftr-bind,lem:liftr-refl}
  \lean{CatCrypt.Relational.rHoare_sample_comm}
  \leanok
  When the left side samples $\alpha$ then $\beta$ and the right side samples
  $\beta$ then $\alpha$, the two orders may be aligned: if
  $\rHoare\,\Phi\,(k_1\,a\,b)\,(k_2\,a\,b)\,\Psi$ for all $a,b$, then the
  differently-ordered sampling programs are related under $\Phi$ and $\Psi$.
\end{theorem}

\section{Advantage-level consequences}

\begin{theorem}[Zero advantage from pRHL equality]
  \label{thm:advantage-zero-of-rhoare}
  \uses{def:rhoare,def:eqpre,def:eqpost,lem:liftr-eq-implies-eq,def:advantageA}
  \lean{CatCrypt.Crypto.advantage_zero_of_rHoare}
  \leanok
  If $\rHoare\,\eqPre\,G_0\,G_1\,\eqPost$, then no adversary can distinguish the
  two games: for every $A$, $\AdvantageA\,G_0\,G_1\,A = 0$. The equality
  postcondition forces the empty-heap output distributions to coincide.
\end{theorem}

\begin{theorem}[Advantage factorization]
  \label{thm:advantage-factorization}
  \uses{def:advantageA}
  \lean{CatCrypt.Relational.advantage_factorization}
  \leanok
  Post-composing both games with the same $f$ moves $f$ into the distinguisher
  without changing the advantage:
  \[
    \AdvantageA\,(G_1 \bind f)\,(G_2 \bind f)\,A
      = \AdvantageA\,G_1\,G_2\,(\lambda x.\ (f\,x) \bind A).
  \]
  The one-directional ``post-processing cannot increase advantage'' bound is the
  immediate consequence.
\end{theorem}

\begin{theorem}[Factorization with a commuted pure prefix]
  \label{thm:advantage-factorization-comm}
  \uses{def:advantageA,thm:advantage-factorization}
  \lean{CatCrypt.Relational.advantage_factorization_comm}
  \leanok
  When both games share a pure prefix $\mathit{pfx}$, it may be commuted past the
  varying parts and factored into the distinguisher: for $\mathsf{IsPure}\,\mathit{pfx}$,
  \[
    \AdvantageA\,(\mathit{pfx} \bind \lambda p.\ G_1 \bind \lambda r.\ k\,p\,r)\,
                 (\mathit{pfx} \bind \lambda p.\ G_2 \bind \lambda r.\ k\,p\,r)\,A
    = \AdvantageA\,G_1\,G_2\,(\lambda r.\ \mathit{pfx} \bind \lambda p.\ (k\,p\,r) \bind A).
  \]
\end{theorem}

\begin{theorem}[Symmetry of advantage]
  \label{thm:advantage-symm}
  \uses{def:advantage}
  \lean{CatCrypt.Relational.advantage_symm}
  \leanok
  The distinguishing advantage is symmetric in its two games:
  $\Advantage\,G_1\,G_2 = \Advantage\,G_2\,G_1$.
\end{theorem}
